\section{Mission Briefing}
\label{sec:dustcastingBriefing}

\subsection{The Objective}

Start this chapter by reading the following:

\quoteBox{
	As you are resting after a difficult day, you see a movement in the only mirror in the room. Looking at it you see a familiar reflection of a black woman with a big eyes, glowing white.
}

If the characters have no access to any of the mirrors (e.g. they are in the Singer camp), instead read the following:

\quoteBox{
	Suddenly, you see your shadow whirl, twist, and grow in a wrong direction - towards the light. Soon, it takes a familiar shape of a woman, with two eyes glowing brightly.
}

Regardless of a way \npcSjaAnat{} appears, continue with the following:

\quoteBox{
	\npcSjaAnat{} smiles and addresses you: "So you survived. That is well, but not unexpected. But your task is yet completed. The city is under siege and now it's your time to make a change that will matter. Take the fabrial from the monastery, the one looking like a bracelet. You will need it sooner than you think."
}

Allow the characters to have a conversation with the Unmade. Use this opportunity to convey the task for the PCs: sneak up and use \lootDustcaster{} to annihilate enemy's biggest asset (the Thunderclast or a barricaded gate).

\scriptedOption{What is this fabrial?}{
	It is \lootDustcaster{}, an ancient gift from \npcDustmother{}, the \npcDustmotherTitle{}.
}

\scriptedOption{And the other one?}{
	They call it \lootGateOpener{}, and it came from \npcBlackFisher{}, the \npcBlackFisherTitle{}, but you shall not need it now. Later, weeks later, yes.
}

\scriptedOption{What does the fabrial do?}{
	It can be used to turn your foes to dust, but only once. Do not use it before the time is right.
}

\scriptedOption{How do we use it?}{
	Adorn your arm with it, and press your palm to your foe. The fabrial will do the rest, it yearns to be awaken again. But you must keep it under control, or it will run free and wild.
}

\scriptedOption{What are the risks of using it?}{
	There are none, if your will reigns over it's desire of destruction. If you fail, it may unmake you as well.
}

\scriptedOption{What should we destroy?}{
	Odium forces just awaken a Thunderclast, a giant old as the war, strong enough to break the walls of the city. But then, the gates are being barricaded and the walls are defended by \npcHighmarshal{} who is stronger than you can imagine. If both sides clash in battle, the river of blood will be deep, and the slain will be many. You must prevent it, tip the scales of the battle, unmake the biggest thread.
}

\subsection{Escape Route}

Once the objective is clear for the PCs, you can add the following:

\quoteBox{
	"Once the deed is done, your enemies will mark you for death and hunt you down", \npcSjaAnat{} continues. "You will need to flee the city. You must enter the Circle of Memories on the Monastery Dais. My son will help you disappear from the enemy eyes. You will need to travel south, until you reach Kharbranth. Use \lootGateOpener{} there, and \npcTaravangian{} will offer you help and shelter. \\ \\ 
	"However, you must first determine how will you reach the Circle of Memories here in Kholinar, as it is now in the hands of the ardents and I can sense Odium influence touching them."
}

Unknown to the PCs, \npcSjaAnat{} suggests their escape rout through a corrupted Oathgate that will transport them to Shadesmar. Previously collected \lootGateOpener{} will help them get back to the Physical Realm from there. 

\reasonBox{
	When you present the escape root to the PCs, don't reveal all the secrets just yet. People who read the books may connect the dots, but for the others there might still be some surprises left. Don't mention Oathgates, Shadesmar and the purpose of \lootGateOpener{} at this stage.
}

However, as the Oathgate is now occupied by \npcFractionCult{}, the characters will need to hatch a plan how to get there. There are several options for that:

\scriptedOption{Brute force.}{
	For sure there will be guards, but the PCs can choose to just kill anyone blocking their path.
}

\scriptedOption{Use connections.}{
	If the party helped enough \enemyCultist{}s so far, they might be friendly enough to just let them in.
}

\scriptedOption{Join the cult.}{
	If the characters prove they are worthy, they might be invited to join the cult. In that case, nobody will try to stop them on their way to the Circle of Memories. \\
}

If the players have difficulty to come up with any idea how to enter the Monastery Dais, you can suggest them a single solution from the above list after passing a \skillCheck{15}{\skillDeduction{}}. Each subsequent suggestion raises difficulty of such test by 5.

\textbf{Note:} Details concerning the escape solutions are described in \fullref{sec:cityEscapePlan}. Despite being located in the future chapter, the PCs should make preparations before they start \fullref{sec:dustcastingSneaking}.

\corruptedBox{
	Meeting or talking to \npcSjaAnat{}, the \npcSjaAnatTitle{}, can expedite the progression of an Enlightened Lightweaver \seeCorruptedRules{}.
}

\corruptedBox{
	Wearing \lootDustcaster{}, an artifact of \npcDustmother{}, the \npcDustmotherTitle{}, for the first time can expedite the progression of an Enlightened Dustbringer \seeCorruptedRules{}.
}



\subsection{Radiant Required}

To activate the Oathgate in the Circle of Memories, the characters will need a living Shardblade summoned by a Radiant Knight. If at this point the party does not include a Radiant, \npcSjaAnat{} will urge the PCs to bond one of her children.
